\hspace{3mm}

\centerline{\bf \Huge Вступительный экзамен. Теория}


\begin{flushright}
        \scriptsize{\textbf{Математика — это единственный совершенный метод водить самого себя за нос. Альберт Эйнштейн}}
\end{flushright}

\problem{\textbf{1}} Напишите определение аксиомы.

\problem{\textbf{2}} Напишите определение делимости двух чисел. 

\problem{\textbf{3}} Напишите определения простых и составных чисел.

\problem{\textbf{4}} Напишите определение ОТА.

\problem{\textbf{5}} Напишите определения НОДа и НОКа.

\problem{\textbf{6}} Единица - это простое или составное число? Напишите развёрнутое объяснение.

\problem{\textbf{7}} Напишите определение взаимнопростых чисел.

\problem{\textbf{8}} Напишите определение диофантовых уравнений?

\problem{\textbf{9}} Нарисуйте все множества чисел с помощью кругов Эйлера.

\problem{\textbf{10}} Напишите определения мощности множества и равномощных множеств.

\problem{\textbf{11}} Напишите определения биекции и счётного множества.

\problem{\textbf{12}} Напишите определение функции.

\problem{\textbf{13}} Чего достаточно для построения прямой и почему этого достаточно?

\problem{\textbf{14}} Как проверить, что точка с определёнными координатами принадлежит графику определённой функции?

\problem{\textbf{15}} Напишите определения пересечения множеств и объединения множеств.

\problem{\textbf{16}} Что такое треугольник Паскаля и для чего он нужен?

\problem{\textbf{17}} Напишите продолжения: Если дискриминант больше нуля, то ...

Если дискриминант равен нулю, то ...

Если дискриминант меньше нуля, то ...

\problem{\textbf{18}} Напишите определение мнимой единицы. С каким множеством она связана?

\begin{flushright} \textbf{2 балла}\\ 
\end{flushright}

\newpage


\problem{\textbf{19}} Верно ли, что все натуральные числа, кроме точных квадратов имеют чётное количество делителей, а точные квадраты - нечётное количество делителей?

Верно ли, что если натуральное число имеет чётное количество делителей, то это не точный квадрат, а если нечётное - то точный квадрат?

\problem{\textbf{20}} Напишите определения ВСЕХ базовых признаков делимости. 

\problem{\textbf{21}} Как можно получить новые признаки делимости? Напишите развёрнутое объяснение и приведите 3 примера новых признаков делимости.

\problem{\textbf{22}} Напишите определения ВСЕХ множеств чисел и ответьте:

Существует ли самое большое натуральное число? Приведите пример или докажите несуществование.

\problem{\textbf{23}} Докажите, что число $\sqrt{2025}$ - рациональное, а $\sqrt{2026}$ - иррациональное. 

\problem{\textbf{24}} Докажите, что множества целых и рациональных чисел счётны.

И ответьте, каких чисел больше: натуральных, которые делятся на 25, или натуральных, которые делятся на 4? Напишите развёрнутое объяснение.

\problem{\textbf{25}} Где больше точек: на отрезке длиной 1 см или на отрезке длиной 7 см? 

Где больше точек: на отрезке длиной 7 см или на бесконечной прямой? 

Где больше точек: на окружности с радиусом 2026 км или на бесконечной прямой?

Напишите развёрнутое объяснение.

\begin{flushright} \textbf{6 баллов}\\ 
\end{flushright}

\problem{\textbf{67}}\textbf{([0;+$\infty))$} Если Вы ничего больше не можете решить и написать, то в этом задании Вы можете написать все теоремы, признаки, свойства, интересные факты и т.д., которые связаны с математикой. 

\centerline{\textbf{Теория: 78+ баллов}}
\centerline{\textbf{Практика: 237 баллов}}
\centerline{\textbf{ИТОГО: 315+ баллов}}


